\subsection{FREE GROUPS}

Let \(G_x = Z\) for all \(x \in X\) . The free product of the family \((G_x)_{x \in X}\) is called the free group constructed on \(X\) and is denoted by \(F(X)\) . Let \(\phi_x\) denote the canonical homomorphism of \(G_x = Z\) into \(F(X)\) . By no. 3, Corollary to Proposition 5, the mapping \(x \mapsto \phi_x(1)\) of \(X\) into \(F(X)\) is injective; we shall identify \(X\) with its image in \(F(X)\) under this mapping. Then \(X\) generates \(F(X)\) and \(e \notin X\) .

Applying no. 3, Proposition 5, we obtain the following result:

PROPOSITION 7. Let \(g\) be an element of the free group \(F(X)\) . There exist an integer \(n \geqslant 0\) and a sequence \((x_{\alpha}, m(\alpha))_{1 \leqslant \alpha \leqslant n}\) determined uniquely by the relations \(x_{\alpha} \in X\) ,

\[
g = \prod_ {\alpha = 1} ^ {n} x _ {\alpha} ^ {m (\alpha)}
\]

The free group \(F(X)\) enjoys the following universal property:

PROPOSITION 8. Let G be a group and f a mapping of X into G. There exists one and only one homomorphism \(\bar{f}\) of F(X) into G which extends \(f\) .

The uniqueness of \(\bar{f}\) follows from the fact that the group \(\mathrm{F}(\mathbf{X})\) is generated by \(\mathbf{X}\) . For all \(x\) in \(\mathbf{X}\) , let \(f_x\) be the homomorphism \(n \mapsto f(x)^n\) of \(\mathbf{Z}\) into \(\mathbf{G}\) . By no. 3, Proposition 4, there exists a homomorphism \(\bar{f}\) of \(\mathrm{F}(\mathbf{X})\) into \(\mathbf{G}\) such that \(\bar{f}(x^n) = f_x(n)\) for \(x \in \mathbf{X}\) and \(n \in \mathbf{Z}\) ; in particular, \(\bar{f}(x) = f_x(1) = f(x)\) for all \(x \in \mathbf{X}\) and hence \(\bar{f}\) extends \(f\) .

Let \(u: X \to Y\) be a mapping. By Proposition 8 there exists one and only one homomorphism of \(F(X)\) into \(F(Y)\) which coincides with \(u\) on \(X\) ; it is denoted by \(F(u)\) . As in the case of magmas (no. 1) the formula

\[
\mathbf {F} (v \circ u) = \mathbf {F} (v) \circ \mathbf {F} (u)
\]

is established for every mapping \(v: \mathbf{Y} \to \mathbf{Z}\) and it is shown that \(\mathrm{F}(u)\) is injective (resp. surjective, bijective) if \(u\) is. For every subset \(S\) of \(X\) , \(F(S)\) will be identified with the subgroup of \(F(X)\) generated by \(S\) .

Let I be a set. In certain cases it is of interest not to identify i in I with its canonical image \(\phi_i(1)\) in the free group \(\mathbf{F}(\mathbf{I})\) ; the latter will be denoted by \(\mathbf{T}_i\) (or \(\mathbf{T}_i', \mathbf{X}_i, \ldots\) as the case may be) and called the indeterminate of index \(i\) . The free group \(\mathbf{F}(\mathbf{I})\) is then denoted by \(\mathbf{F}((\mathbf{T}_i)_{i \in 1})\) or \(\mathbf{F}(\mathbf{T}_1, \ldots, \mathbf{T}_n)\) if \(\mathbf{I} = \{1, 2, \ldots, n\}\) .

Let G be a group and \(\mathbf{t} = (t_{i})_{i \in \mathbf{I}}\) a family of elements of G. By Proposition 8 there exists a homomorphism \(f_{t}\) of \(\mathrm{F}((\mathrm{T}_{i})_{i \in \mathbf{I}})\) into G characterized by \(f_{\mathbf{t}}(\mathrm{T}_{i}) = t_{i}\) for all \(i \in I\) . The image of an element w of \(\mathrm{F}((\mathrm{T}_{i})_{i \in \mathbf{I}})\) under \(f_{t}\) will be denoted by \(w(\mathbf{t})\) or \(w(t_{1}, \ldots, t_{n})\) if \(I = \{1, 2, \ldots, n\}\) ; \(w(\mathbf{t})\) is said to result from the substitution \(T_{i} \mapsto t_{i}\) in w. In particular, if we take \(\mathrm{G} = \mathrm{F}((\mathrm{T}_{i})_{i \in \mathbf{I}})\) and \((t_{i}) = (\mathrm{T}_{i}) = \mathbf{T}\) , \(f_{T}\) is the identity homomorphism of G, whence \(w(\mathbf{T}) = w\) ; for \(I = \{1, 2, \ldots, n\}\) , then \(w(\mathrm{T}_{1}, \ldots, \mathrm{T}_{n}) = w\) .

Let G and G' be two groups, u a homomorphism of G into G' and \(\mathbf{t} = (t_{1}, \ldots, t_{n})\) a finite sequence of elements in G. Let \(\mathbf{t}' = (u(t_{1}), \ldots, u(t_{n}))\) ; the homomorphism \(u \circ f_{t}\) of \(\mathrm{F}(\mathrm{T}_{1}, \ldots, \mathrm{T}_{n})\) into G' maps \(T_{i}\) to \(u(t_{i})\) for \(1 \leqslant i \leqslant n\) and hence is equal to \(f_{t'}\) ; for w in \(\mathrm{F}(\mathrm{T}_{1}, \ldots, \mathrm{T}_{n})\) , then

\[
u (w (t _ {1}, \dots , t _ {n})) = w (u (t _ {1}), \dots , u (t _ {n})).\tag{8}
\]

Let \(w\) be given in \(F(T_1, \ldots, T_n)\) and elements \(v_1, \ldots, v_n\) in the free group \(F(T_1', \ldots, T_m')\) . The substitution \(T_i \mapsto v_i\) defines an element \(w' = w(v_1, \ldots, v_n)\) of \(F(T_1', \ldots, T_m')\) . Let \(G\) be a group, \(t_1, \ldots, t_m\) elements of \(G\) and \(u\) the homomorphism of \(F(T_1', \ldots, T_m')\) into \(G\) characterized by \(u(T_j') = t_j\) for \(1 \leqslant j \leqslant m\) . Then \(u(v_i) = v_i(t_1, \ldots, t_m)\) and \(u(w') = w(t_1, \ldots, t_m)\) ; formula (8) thus implies

\[
w ^ {\prime} (t _ {1}, \dots , t _ {m}) = w (v _ {1} (t _ {1}, \dots , t _ {m}), \dots , v _ {n} (t _ {1}, \dots , t _ {m})).\tag{9}
\]

This justifies the ``functional notation'' \(w(t_{1},\ldots,t_{n})\) . The reader is left to extend formulae (8) and (9) to the case of arbitrary indexing sets.

